
\documentclass{article}
\usepackage{colortbl}
\usepackage{makecell}
\usepackage{multirow}
\usepackage{supertabular}

\begin{document}

\newcounter{utterance}

\centering \large Interaction Transcript for game `eqbench', experiment `eq\_bench', episode 61 with reflexion{-}Qwen3.5{-}27B.
\vspace{24pt}

{ \footnotesize  \setcounter{utterance}{1}
\setlength{\tabcolsep}{0pt}
\begin{supertabular}{c@{$\;$}|p{.15\linewidth}@{}p{.15\linewidth}p{.15\linewidth}p{.15\linewidth}p{.15\linewidth}p{.15\linewidth}}
    \# & \multicolumn{2}{c}{Player} && \multicolumn{2}{c}{Game Master} \\
    \hline

    \theutterance \stepcounter{utterance}  
    & & & \multicolumn{4}{p{0.6\linewidth}}{
        \cellcolor[rgb]{0.9,0.9,0.9}{
            \makecell[{{p{\linewidth}}}]{
                \texttt{\tiny{[P1$\langle$GM]}}
                \texttt{You have attempted this same task before and did not succeed. These are the plans you wrote after those attempts.} \\
\texttt{Trial \#1: New plan: My previous attempt failed because I assigned a high score to "Intrigue" (7) despite the dialogue's clear trajectory toward a direct confrontation and threat. Julia's responses ("I'm not afraid," "Is that a threat?") explicitly reject fear and challenge Marcus, making "Defiance" the dominant emotion. "Intrigue" is a secondary reaction to the mystery, but the immediate emotional response to a "warning" and "threat" is resistance, not curiosity. Furthermore, "Indifference" should be 0 as she is actively engaged, and "Fear" should be very low (0 or 1) as she explicitly states she is not afraid. Next time, I will prioritize the explicit emotional stance of the character (defiance against a threat) over the thematic context (mystery), ensuring the highest score reflects the most immediate and intense reaction to the final lines. I will adjust the scores to: Fear: 0, Defiance: 10, Indifference: 0, Intrigue: 2.} \\
\\ 
\\ 
\\ 
\texttt{Your task is to predict the likely emotional responses of a character in this dialogue:} \\
\\ 
\texttt{Marcus: Stuck in an elevator. Just my luck.} \\
\texttt{Julia: Could be worse. Could be stuck in a room with no windows.} \\
\texttt{Marcus: You're an optimist.} \\
\texttt{Julia: Realist. I just choose to see the glass half full.} \\
\texttt{Marcus: And what if the glass is poisoned?} \\
\texttt{Julia: Then it's a good thing I'm thirsty.} \\
\texttt{Marcus: You always this reckless?} \\
\texttt{Julia: Only when it matters.} \\
\texttt{Marcus: And what matters to you?} \\
\texttt{Julia: The truth.} \\
\texttt{Marcus: Truth is a dangerous game.} \\
\texttt{Julia: Only for those with something to hide.} \\
\texttt{Marcus: Or those who seek it.} \\
\texttt{Julia: I'm not afraid.} \\
\texttt{Marcus: You should be.} \\
\texttt{Julia: Is that a threat?} \\
\texttt{Marcus: A warning.} \\
\texttt{{[}End dialogue{]}} \\
\\ 
\texttt{At the end of this dialogue, Julia would feel...} \\
\texttt{Fear} \\
\texttt{Defiance} \\
\texttt{Indifference} \\
\texttt{Intrigue} \\
\\ 
\texttt{Give each of these possible emotions a score from 0{-}10 for the relative intensity that they are likely to be feeling each.} \\
\\ 
\texttt{You must output in the following format, including headings (of course, you should give your own scores), with no additional commentary:} \\
\\ 
\texttt{Fear: <score>} \\
\texttt{Defiance: <score>} \\
\texttt{Indifference: <score>} \\
\texttt{Intrigue: <score>} \\
\\ 
\\ 
\texttt{{[}End of answer{]}} \\
\\ 
\texttt{Remember: zero is a valid score, meaning they are likely not feeling that emotion. You must score at least one emotion > 0.} \\
\\ 
\texttt{Your answer:} \\
            }
        }
    }
     \\ \\

    \theutterance \stepcounter{utterance}  
    & \multicolumn{4}{p{0.6\linewidth}}{
        \cellcolor[rgb]{0.561,0.737,0.561}{
            \makecell[{{p{\linewidth}}}]{
                \texttt{\tiny{[P1$\rangle$GM]}}
                \texttt{Fear: 0} \\
\texttt{Defiance: 10} \\
\texttt{Indifference: 0} \\
\texttt{Intrigue: 2} \\
            }
        }
    }
    & & \\ \\

    \theutterance \stepcounter{utterance}  
    & & & \multicolumn{2}{p{0.3\linewidth}}{
        \cellcolor[rgb]{0.9,0.9,0.9}{
            \makecell[{{p{\linewidth}}}]{
                \texttt{\tiny{[GM$|$GM]}}
                \texttt{Fear: 0} \\
\texttt{Defiance: 10} \\
\texttt{Indifference: 0} \\
\texttt{Intrigue: 2} \\
            }
        }
    }
    & & \\ \\

    \theutterance \stepcounter{utterance}  
    & & & \multicolumn{2}{p{0.3\linewidth}}{
        \cellcolor[rgb]{0.9,0.9,0.9}{
            \makecell[{{p{\linewidth}}}]{
                \texttt{\tiny{[GM$|$GM]}}
                \texttt{Fear: 3} \\
\texttt{Defiance: 7} \\
\texttt{Indifference: 0} \\
\texttt{Intrigue: 6} \\
            }
        }
    }
    & & \\ \\

    \theutterance \stepcounter{utterance}  
    & & & \multicolumn{2}{p{0.3\linewidth}}{
        \cellcolor[rgb]{0.9,0.9,0.9}{
            \makecell[{{p{\linewidth}}}]{
                \texttt{\tiny{[GM$|$GM]}}
                \texttt{game\_result = LOSE} \\
            }
        }
    }
    & & \\ \\

    \theutterance \stepcounter{utterance}  
    & & \multicolumn{4}{p{0.6\linewidth}}{
        \cellcolor[rgb]{0.95,0.95,0.95}{
            \makecell[{{p{\linewidth}}}]{
                \texttt{\tiny{[P1$\leftrightarrow$P]}}
                \texttt{Post{-}episode Reflexion (not sent to GM):} \\
\texttt{New plan: My previous attempt failed because I assigned a non{-}zero score to "Intrigue" (2), which diluted the primary emotional response. The dialogue culminates in a direct threat ("A warning") to which Julia's established character trait is explicit fearlessness ("I'm not afraid") and active resistance ("Is that a threat?"). The prompt asks for the *likely* emotional response at the very end; in the face of a direct warning, her reaction is pure confrontation, not curiosity. To achieve the target state, I must assign the maximum score (10) exclusively to "Defiance" and set all other emotions, including "Intrigue," to 0. This ensures the output reflects the singular, dominant emotional state of a character who refuses to be intimidated. I will output: Fear: 0, Defiance: 10, Indifference: 0, Intrigue: 0.} \\
            }
        }
    }
    & \\ \\

\end{supertabular}
}

\end{document}
